\chapter{Introduction to LLM Agents}
\label{ch:introduction}

\epigraph{``An agent is not a thing; it is a pattern. Any system that perceives, decides, and acts is an agent.''}{}

\learninggoal{
	\item Define what an LLM agent is and distinguish it from a bare LLM.
	\item Trace the intellectual lineage from classical planning to LLM-based agents.
	\item Identify the core components that every agent system must implement.
	\item Understand when an agent architecture is needed versus a direct prompt.
}

\section{What Is an Agent?}

An \emph{agent} is a system that perceives its environment, makes decisions, and takes actions to achieve goals. In the context of large language models, an agent wraps an \llm\ with three additional capabilities:

\begin{enumerate}
	\item \textbf{Structured reasoning:} the ability to decompose problems, evaluate alternatives, and reflect on outcomes.
	\item \textbf{Tool use:} the ability to call external functions (APIs, databases, code executors, search engines) and incorporate their results.
	\item \textbf{State management:} the ability to maintain memory across interactions, including short-term (context window), long-term (persistent storage), and episodic memory.
\end{enumerate}

\begin{definition}[LLM Agent]
	An \textbf{LLM agent} is a software system in which a language model serves as the core reasoning engine, augmented with:
	\begin{itemize}
		\item A \emph{perception module} that converts environment observations into text.
		\item An \emph{action module} that maps LLM outputs to environment interventions.
		\item A \emph{memory module} that persists state across reasoning steps.
		\item A \emph{tool registry} that describes available functions and their invocation protocols.
	\end{itemize}
\end{definition}

\keyinsight{The defining property of an agent is that its outputs have \emph{side effects}. A bare LLM generates text; an agent generates text that changes the world---by writing files, sending emails, executing code, or controlling robots.}

\section{Intellectual Lineage}

LLM agents did not emerge from a vacuum. They synthesise three research traditions:

\subsection{Classical Planning (1956--1990s)}

STRIPS~\cite{wohlman1977strips}, SOAR~\cite{laird1987soar}, and ACT-R~\cite{anderson1996actr} modelled agents as symbolic reasoners with explicit world models. A planner maintained a state representation, searched over action sequences, and executed the optimal plan. The limitations:
\begin{itemize}
	\item Brittle: plans broke when the world deviated from the model.
	\item Knowledge bottleneck: encoding common sense into symbolic rules proved impractical.
	\item Scaling: search over action spaces grew exponentially.
\end{itemize}

LLM agents inherit the planning loop (perceive $\to$ reason $\to$ act) but replace the symbolic planner with a neural language model that brings broad world knowledge.

\subsection{Reinforcement Learning Agents (1990s--2020s)}

RL agents learned policies through trial-and-error interaction~\cite{sutton1998rl}. Deep RL (DQN, PPO) extended this to high-dimensional observation spaces. Limitations:
\begin{itemize}
	\item Sample inefficiency: millions of episodes needed for simple tasks.
	\item Reward design: specifying good reward functions is difficult.
	\item Transfer: policies learned in one environment rarely transfer.
\end{itemize}

LLM agents inherit the perception-action loop from RL but replace the learned policy with few-shot-prompted reasoning, drastically reducing the samples needed for new tasks.

\subsection{Augmented Language Models (2020--2023)}

The direct precursor: giving LLMs access to tools. Toolformer~\cite{schick2023toolformer} showed that LLMs could learn to call APIs from a few demonstrations. WebGPT~\cite{nakano2021webgpt} demonstrated web browsing agents. The step from ``LLM with a tool call'' to ``agent'' came when researchers added iterative reasoning, memory, and reflection~\cite{yao2023react}.

\section{When Do You Need an Agent?}

The distinction between a prompted LLM and an agent is not binary---it is a spectrum of autonomy:

\begin{longtable}{p{3cm}p{4cm}p{4cm}p{4cm}}
	\toprule
	\textbf{Dimension} & \textbf{Bare LLM} & \textbf{Augmented LLM}  & \textbf{Full Agent}      \\
	\midrule
	Output             & Text              & Text + structured calls & Multi-step actions       \\
	Context            & Single prompt     & Prompt + retrieved docs & History + state + tools  \\
	State              & Stateless         & Cache-aware             & Persistent memory        \\
	Error recovery     & None              & Retry on parse failure  & Reflection + replan      \\
	Environment        & None              & Read-only (search)      & Read-write (APIs, files) \\
	Autonomy           & None              & Tool suggestions        & Autonomous execution     \\
	\bottomrule
\end{longtable}

\begin{example}[When to Use an Agent]
	\begin{itemize}
		\item \textbf{Direct LLM is fine:} summarise an article, translate a paragraph, answer a factual question.
		\item \textbf{Augmented LLM helps:} answer a question that requires web search, generate an email from a template.
		\item \textbf{Agent required:} book a flight by searching multiple sites, comparing prices, filling forms, and confirming. Write and test a multi-file code change. Manage a supply chain across dozens of vendors.
	\end{itemize}
\end{example}

\section{The Agent Loop}

Every agent system, regardless of complexity, implements a core loop:

\begin{lstlisting}[caption=The universal agent loop]
  state = initial_state()
  while not done:
      observation = perceive(environment)
      state    = update_memory(state, observation)
      action   = reason(state, tools, goal)
      result   = execute(action)
      state    = update_memory(state, result)
      done     = is_complete(state, goal)
\end{lstlisting}

This loop appears in every framework---LangChain, AutoGen, CrewAI, Semantic Kernel---with variations in how \texttt{reason()} is implemented. Later chapters examine each component in depth.

\section{How This Book Is Organised}

This book is divided into three parts:

\begin{description}
	\item[Part I: Core Architecture (Chapters 2--5)] covers the fundamental building blocks: the agent loop, reasoning frameworks, tool use, and memory systems. Every agent system must solve these problems.
	\item[Part II: Advanced Topics (Chapters 6--8)] covers planning, multi-agent systems, and grounding/reliability. These capabilities distinguish hobby projects from production systems.
	\item[Part III: Engineering (Chapters 9--11)] covers evaluation, safety, and production deployment. These chapters are about making agents robust enough for real users.
\end{description}

\section{Recommended Reading}

\begin{itemize}
	\item Xi et al.~\cite{xi2023rise} -- comprehensive survey of LLM-based agents.
	\item Sumers et al.~\cite{sumers2024cognitive} -- cognitive architecture perspective on language agents.
	\item Mialon et al.~\cite{mialon2023augmented} -- survey of augmented language models.
\end{itemize}

\section{Exercises}

\begin{exercise}
	Identify a task you currently do manually (e.g., triaging email, filing expenses, updating a wiki). Sketch the agent loop that could automate it. What tools would it need? What could go wrong?
\end{exercise}

\begin{exercise}
	Read the definition of an agent in~\cite{russell2016aima}. How does the LLM agent definition differ? Where do LLM agents fall short of the classical definition?
\end{exercise}

\begin{exercise}
	Take a simple LLM application (e.g., a chatbot that answers from documentation). Trace the minimum changes needed to turn it into an agent. Where does the complexity come from?
\end{exercise}
